\section{Introduction and conventions}\label{sec:introduction}

Let $\Omega,\Lambda\subset\R^3$ be domains, meaning nonempty connected
open sets.  A homeomorphism $f:\Omega\to\Lambda$ belongs to
$W^{1,p}(\Omega;\R^3)$ when its components and their weak first
derivatives are in $L^p$.  Strong $W^{1,p}$ approximation requires
convergence of both the maps and these derivatives in $L^p$.
The \emph{fixed-target Ball--Evans approximation problem} asks whether
one can choose the approximants to be smooth diffeomorphisms from
$\Omega$ onto exactly $\Lambda$.  Here a diffeomorphism and its inverse
are smooth on their open domains; this definition makes no assertion
about boundary values.

The problem originates in nonlinear elasticity, where injectivity
expresses noninterpenetration and the derivative describes local
deformation.  Ball's discussions~\cite{Ball2001,Ball2002,Ball2010}
connect the approximation question to the variational and computational
theory; the original planar approximation paper records his attribution
of the question to Evans~\cite{IwaniecKovalevOnninen2011}.
Neither ordinary convolution nor sufficiently fine vertex interpolation
automatically preserves injectivity.  We resolve the fixed-target
problem on bounded three-dimensional domains for every finite exponent
above two.

\begin{theorem}\label{main:theorem}\label{ha:approximation}
Let $\Omega,\Lambda\subset\R^3$ be bounded domains, let $2<p<\infty$,
and let $f:\Omega\to\Lambda$ be a homeomorphism in
$W^{1,p}(\Omega;\R^3)$. There exist $C^\infty$ diffeomorphisms
$f_j:\Omega\to\Lambda$, each belonging to
$W^{1,p}(\Omega;\R^3)$, such that
\begin{equation}\label{main:convergence}
\int_\Omega\bigl( |f_j-f|^p+|Df_j-Df|^p\bigr)\,dx\longrightarrow0.
\end{equation}
\end{theorem}

Neither boundary regularity, a boundary extension of $f$, inverse
Sobolev regularity, nor a nonvanishing Jacobian is assumed.  The
theorem concerns forward Sobolev convergence.  Smoothness of each
inverse does not assert Sobolev convergence of the inverses; strong
forward convergence alone also does not imply convergence of an elastic
energy with a singular determinant penalty.  These are distinct
approximation questions; see \cite{Ball2002,Pratelli2017,Hencl2025}.
The complementary range $1\le p\le2$ is treated by a different
construction in the companion article~\cite{OpenAILow2026}.
No theorem from that range is used here.  The full common smoothing
argument is included in Appendix~\ref{sec:smoothing}, so the present
range theorem has a complete local proof.

\subsection{Earlier results and the dimensional obstacle}

The planar theory provides both the principal positive precedents and
important contrasts.  Iwaniec, Kovalev, and Onninen treated the
Hilbert-space case $p=2$ using harmonic replacements in work first
circulated in June 2010~\cite{IwaniecKovalevOnninen2012}.  Their extension,
first circulated that September, covers every $1<p<\infty$ on arbitrary planar open sets
\cite[Theorem~1.1]{IwaniecKovalevOnninen2011}.  Hencl and Pratelli
proved the endpoint $p=1$ by geometric extension and grid constructions,
including locally finite piecewise affine approximation
\cite[Theorem~1.1]{HenclPratelli2018}.  Both results preserve the
target and require no inverse Sobolev hypothesis.  Thus those features
are inherited goals of the problem, not new distinctions of the
three-dimensional statement.

Geometric approximation has also been developed under stronger planar
input assumptions.  Bellido and Mora-Corral control approximation in
smaller H\"older exponents on polygonal planar domains when both
directions are H\"older
\cite{BellidoMoraCorral2011}.  Daneri and Pratelli obtain quantitative
bi-Lipschitz approximation with derivative control in both directions
for bi-Lipschitz maps~\cite{DaneriPratelli2014}; Pratelli treats
bi-$W^{1,1}$ maps~\cite{Pratelli2017}.  These results clarify the
distinction between approximating a map by maps with smooth inverses
and approximating the inverse in a specified Sobolev norm.

Piecewise affine smoothing is a separate stage.  Mora-Corral and
Pratelli established a strong planar version with control in both
directions~\cite{MoraCorralPratelli2014}.  Campbell and Soudsk\'y
obtained $C^1$ injective approximation of piecewise affine maps in
arbitrary dimension, without asserting smoothness of the inverse
\cite{CampbellSoudsky2021}.  Campbell, D'Onofrio, and V\'itek then
proved diffeomorphic approximation of locally finite piecewise affine
homeomorphisms in dimensions three and four, with uniform and
derivative-error control for both the map and its inverse
\cite[Theorem~A]{CampbellDOnofrioVitek2026}.  The remaining task for a
general Sobolev input is to obtain the required strong approximation
before this smoothing stage.  Appendix~\ref{sec:smoothing} includes
both the locally bi-Lipschitz-to-PL reduction and a PL smoothing proof
with arbitrary positive continuous value tolerances.

The dimension restriction is substantive.  Building on the $W^{1,1}$
construction of Hencl and Vejnar~\cite{HenclVejnar2016}, and repairing
a step in that construction, Campbell, Hencl, and Tengvall obtained
counterexamples in dimensions $n\ge4$ for
$1\le p<\lfloor n/2\rfloor$, with Jacobians of both signs on sets
of positive measure~\cite{CampbellHenclTengvall2018}.  Those examples
do not decide the three-dimensional problem.  Qualitative topology
does help in dimension three: the triangulation and PL-approximation
theory of Moise and Bing, in Hamilton's relative formulation,
provides the topological flexibility used below
\cite{Moise1952Approximation,Moise1952Triangulation,Bing1959,Hamilton1976}.
It supplies no strong derivative estimate.  In particular, a
topological extension chosen after a small-energy exceptional set may
have an arbitrarily large Lipschitz constant, so multiplying that
earlier energy by its later constant is not a valid approximation
argument.

\subsection{Proof strategy and organization}

We first seek a locally bi-Lipschitz homeomorphism strongly close to
$f$.  At points where $Df$ is invertible, suitably chosen small
cubes permit exact affine replacements.  The principal work concerns
the regular rank-one and rank-two data.  Their image is volume-null,
but their source Sobolev energy need not be small.  Discarding them
as an exceptional set would therefore lose the approximation theorem.
The construction retains them and recovers their gradients through
two scalar coordinates.

On each small retained patch, these coordinates have prescribed affine
first-order models.  Keeping the realized labels close to the original
scalar coordinates controls the change in their mean gradients through
a boundary integral.  Nearly
sharp upper bounds for the gradient norms then give strong gradient
closeness by uniform convexity.  This is why the construction must
control both label values and scalar parameter slopes.

\begin{enumerate}
\item \emph{Prepare exact regions and retained deficient-rank data.}
Section~\ref{sec:high-preparation} first uses the decomposability bundle
of Alberti--Marchese and the converse to Rademacher's theorem of
De Philippis--Rindler to identify transverse directions in the transported
measure~\cite{AlbertiMarchese2016,DePhilippisRindler2016}.
Thin graph bands, related to the null-set constructions of
Alberti--Cs\"ornyei--Preiss~\cite{AlbertiCsornyeiPreiss2010}, flatten
the retained data onto separated plate patches while preserving
their tangential derivatives.  These patches determine the exact
regions required of a continuous proper surjection
$F_b:\Omega\to\Lambda$, called the carrier.  It agrees with affine
or PL homeomorphisms over a prescribed open part of the target and
has controlled local Sobolev energy outside the marked exceptional
interiors that may occur when $2<p<3$.  Its only non-singleton fibres
are specified tame balls, and the proof gives their relative opening
to homeomorphisms.  The cubical cutoff uses immersions of a
two-dimensional skeleton, with the positive-codimension immersion
theorem of Hirsch~\cite{Hirsch1959} and locally flat
Schoenflies~\cite{Brown1960}.

\item \emph{Reduce the variable target to a two-dimensional complex.}
Section~\ref{sec:quantizers} constructs a Lipschitz map
$T:\Lambda\to\Lambda$ from a strongly convex potential.
Its full-dimensional input blocks lie entirely in the carrier's exact
regions; outside these blocks its image is a polyhedral two-complex.
This combines classical convex conjugacy and polyhedral diagram
geometry~\cite{Moreau1965,Aurenhammer1987}, with additional protected
paired contacts proved here.  They place the retained deficient-rank
data on rectangular faces with nearly unchanged gradients.
On a polygonal face, the two coordinates are a scaled logarithm of
the normalized radial fraction and arclength around the perimeter.

\item \emph{Realize these coordinates by source walls.}
Section~\ref{sec:walls} uses annular walls for radial levels and
disk walls for perimeter levels.  Routing removes unwanted
circle intersections without accumulating a derivative factor at
each switch or losing the length of a wall's transverse parameter
interval.  A family of reference fibre paths, called guards,
then pins the realized labels: a misplaced wall would force a
fixed positive path cost on too many small spheres.  The resulting
lower bound grows as the mesh scale tends to zero because $p>2$.
Theorem~\ref{wl:realization} is explicitly conditional on the
geometric and energy estimates proved in
Sections~\ref{sec:mesh}--\ref{sec:high-assembly}.

\item \emph{Supply calibrated meshes and parameter bounds.}
Section~\ref{sec:mesh} prepares and normalizes the source walls,
preserving the individual edge-count bounds, and constructs their
two-colour neighborhoods with simultaneous product coordinates.
Normalization and cut-and-paste belong to the normal-surface
tradition~\cite{Haken1961}.  Section~\ref{sec:calibrated-islands}
improves the estimates on selected source patches.  In aligned source
coordinates, the two scalar gradients are nearly orthogonal in rank two and
nearly parallel in rank one.  The snapping map in
Lemma~\ref{ms:snapping} removes the large parameter slopes caused
by thin intermediate configurations.  Its uniform estimate is only
in the direction from new to old coordinates, as illustrated in
Figure~\ref{ms:fig-snapping}.  Broadening then gives the sharp
scalar bounds by direct differentiation, not by differentiating an
uncontrolled ambient extension.

\item \emph{Recover strong gradients and smooth with a fixed target.}
Section~\ref{sec:high-assembly} chooses local upper bounds for
$\int G^p$, where $G$ is the wall-parameter slope density, before
choosing guards and the final mesh.  These are the capacities used
in the wall argument, not Sobolev capacities of sets.  Nearly sharp
energy and pinned mean gradients yield strong gradient approximation
by uniform convexity.  A separate compression handles marked
exceptional interiors only when $2<p<3$; for $p\ge3$ the volume
Lusin property eliminates them.  The proof ends by correcting the
reserved full-rank cubes and applying the complete smoothing theorem
of Appendix~\ref{sec:smoothing}.  The strict target maps are
self-homeomorphisms, and the final local replacements preserve their
old images; the resulting homeomorphic approximants are onto the
fixed target.
\end{enumerate}

The central quantitative distinction is between derivatives of the
two retained parameters and derivatives of a homeomorphism filling
the complementary chambers.  Only the former enter the principal
energy estimates.  The constant and error choices are recorded in
Remark~\ref{ha:scale-order}; they ensure that every uncontrolled
constant is fixed before the error required to absorb it.

\subsection{Conventions}

Unless stated otherwise, all domains, closures, local finiteness, and compactness are taken relative to the indicated open manifold. A compact set in an open domain has positive distance from its Euclidean complement. A PL map or triangulation is locally finite. A locally bi-Lipschitz homeomorphism has finite Lipschitz constants in both directions on sufficiently small neighborhoods; no global bound is implicit.

The symbol $|A|$ denotes the Frobenius norm of a matrix, and $\|A\|_{\op}$ its operator norm. We use the Frobenius norm in sharp gradient estimates. Equivalent norms give the same strong convergence in Theorem~\ref{main:theorem}. The symbols $\cL^m$ and $\cH^m$ denote Lebesgue and Hausdorff measure. Integrals on parameterized traces include their stated multiplicities. The differential along an $m$-dimensional parameter space is denoted $D_m$ when it must be distinguished from the full differential.

A \emph{fine value tolerance} on a domain $U$ is a positive continuous function on $U$. To approximate finely means to meet a prescribed such tolerance pointwise. Whenever inverse control is needed on a compact localization, it is included in the prescribed tests. Locally finite constructions permit tolerances tending to zero toward the ends of the open domains; no positive global lower bound is assumed.

Constants called \emph{absolute} depend only on dimension and on fixed universal reference shapes. A constant $C_p$ may also depend on the fixed exponent. Other dependencies are specified when the constant is introduced. A finite constant that depends on a chosen compact collection of jets or charts is fixed before an error is required to be small against it.

A homeomorphism between connected oriented domains has a constant topological orientation. If necessary, reflect the target, perform the construction for the orientation-preserving map, and undo the reflection. This convention imposes no extra hypothesis on the Sobolev differential.

\begin{lemma}[Local tolerances]\label{pre:local-tolerances}
Let $U\subset\R^n$ be open and let $a:U\to(0,\infty)$ be continuous. For every $L>0$ there is a positive $L$-Lipschitz function $d$ on $U$ such that
\[
 d(x)\le \min\{a(x),1,L\dist(x,\R^n\setminus U)\}.
\]
There are also positive smooth minorants satisfying prescribed positive continuous upper bounds on their values and first derivatives. One may additionally impose positive constant bounds on a locally finite family of compact regional tests. Finite collections of such requirements can be combined, and error budgets on a countable compact exhaustion can be made summable.
\end{lemma}
\begin{proof}
Extend $e(x)=\min\{a(x),1,L\dist(x,\R^n\setminus U)\}$ by zero outside $U$, and put
\[
 d(x)=\inf_{y\in\R^n}\{e(y)+L|x-y|\}.
\]
This function is $L$-Lipschitz and is bounded above by $e$. It is positive at an interior point $x$: on a small closed ball about $x$, $e$ has a positive minimum, and outside that ball the distance term has a positive lower bound.

For a smooth minorant take a smooth locally finite partition of unity $\{\chi_i\}$ with compact supports in $U$ and put $b=\sum_i c_i\chi_i$, with positive constants $c_i$. Given positive continuous bounds $v,w$, choose
\[
 c_i\le 2^{-i-2}\min\left\{
 \inf_{\supp\chi_i}v,
 \frac{\inf_{\supp\chi_i}w}{1+\|D\chi_i\|_\infty}
 \right\}.
\]
Then $b>0$, $b\le v$, and $|Db|\le w$. Replace a finite list of requirements by their minimum. For countably many regional tests use a locally finite compact cover with compact enlargements, refine the bounds on each intersecting support, and allocate total errors by a convergent positive series. Only finitely many regional requirements meet a given compact support.
\end{proof}

We will also use the following elementary global observation. If $H:U\to V$ is a homeomorphism and a continuous map $G:U\to V$ satisfies
\[
 |G(x)-H(x)|<\tfrac12\dist(H(x),\R^n\setminus V),
\]
then the straight homotopy stays in $V$. When $V$ is bounded, it is a proper homotopy: the preimage of a compact target set is contained, uniformly in the homotopy parameter, in the $H$-preimage of a compact set a positive distance from $\partial V$. This observation will be used together with local injectivity and degree, not as a substitute for local injectivity.
